
Or
\[
\begin{aligned}
  f_*(\underline{F}(n))
  &=\bigoplus f_*\bigl(f^*(\underline{E}_j)\otimes_{\underline{O}_X}\underline{O}_X(n+j)\bigr)\\
  &=\bigoplus \underline{E}_j\otimes_{\underline{O}_S} f_*(\underline{O}_X(n+j)).
\end{aligned}
\]

Par
ailleurs, \(f_*(\underline{O}_X(n+j))=\Sym(\underline{E})^{n+j}\)~; nous noterons \(\sigma_p\) la classe dans \(K^{\bullet}(S)\) de \(\Sym(\underline{E})^p\). La relation (1.11.1) donne donc finalement :
\begin{equation}\tag{1.11.2}
  \forall n\geq n_0,
  \qquad
  \sum_{j=0}^r x_j\sigma_{n+j}=0.
\end{equation}

On considère alors le complexe :
\[
  0\longrightarrow \bigwedge^{r+1}(\underline{E})\otimes_{\underline{O}_S}\Sym(\underline{E})
  \longrightarrow \bigwedge^r(\underline{E})\otimes_{\underline{O}_S}\Sym(\underline{E})
  \longrightarrow\ldots\longrightarrow
  \underline{E}\otimes_{\underline{O}_S}\Sym(\underline{E})
  \longrightarrow \Sym(\underline{E})\longrightarrow0,
\]
dont les flèches sont définies par :
\[
\begin{aligned}
&(x_1\wedge\ldots\wedge x_n)\otimes(y_1\otimes\ldots\otimes y_p)\rightsquigarrow{}\\
&\qquad \sum(-1)^{i-1}(x_1\wedge\ldots\wedge\widehat{x_i}\wedge\ldots\wedge x_n)
\otimes(x_i\otimes y_1\otimes\ldots\otimes y_p)
\end{aligned}
\]
Si on considère \(\bigwedge^n(\underline{E})\) comme \(\underline{O}_S\)-Module gradué réduit au degré \(n\), c'est un complexe de \(\Sym(\underline{E})\)-Modules gradués, les morphismes étant eux-mêmes gradués.

Ce complexe est une résolution de \(\underline{O}_S\) : sur un ouvert où \(\underline{E}\) est libre, \(\Sym(\underline{E})\) s'identifie à une Algèbre de polynomes \(\underline{O}_S[T_0,\ldots,T_r]\), et le complexe au complexe de Koszul défini par la suite \(\Sym(\underline{E})\)-régulière \(T_0,\ldots,T_r\); on sait que ce dernier est une résolution de \(\Sym(\underline{E})/(T_0,\ldots,T_r)=\underline{O}_S\).

En écrivant que chacune des composantes homogènes du complexe de degré \(\geq1\) est acyclique, on obtient donc des relations dans \(K^{\bullet}(S)\) :
\begin{equation}\tag{1.11.3}
  \sum_{i=0}^{r+1}(-1)^i\lambda^i(E)\sigma_{\rho-i}=0
  \qquad\text{pour tout }p\geq1.
\end{equation}
(en posant \(\sigma_k=0\) pour \(k<0\)).

Supposons la relation (1.11.2) vraie pour \(n\geq n_0\); nous allons montrer qu'elle est encore vraie pour \(n_0-1\), ce qui permettra de redescendre par récurrence jusqu'à \(n_0=-r\) (nous verrons plus loin que cela suffit à entrainer la nullité des \(x_j\)).

